\section{Reproducibility}
\label{app:repro}

\paragraph{Environment.}
\texttt{scripts/kaggle\_setup.sh} installs LeRobot v0.6.1 with the \texttt{smolvla} and \texttt{libero} extras, pins MuJoCo below 3.4, registers the NVIDIA EGL vendor library for headless rendering, and answers LIBERO's first-import prompt.
Versions: Python 3.13, PyTorch 2.11 (CUDA 12.8), transformers 5.5.4, hf-libero 0.1.4, MuJoCo 3.3.7.

\paragraph{Evaluation.}
\texttt{python -m evla.run\_eval} runs one configuration on a list of (task, initial state) episodes and appends one CSV row per episode, so that interrupted runs resume and runs from several sessions can be concatenated.
Each episode uses the environment seed $1000+e$ with $e=1000\cdot\text{task}+\text{init}$, and the flow noise of the re-plan at control step $k$ is drawn from a generator seeded with $10^4 e+k$ for every method (common random numbers).
\texttt{scripts/run\_queue.sh} runs a list of configurations as four processes on the two GPUs and pushes partial results every 20 minutes.
\texttt{python -m evla.analyze} produces the paired comparisons, \texttt{python -m evla.bench\_graph} the latency table, and \texttt{python -m evla.profile} the component breakdown.

\paragraph{Variants.}
\cref{tab:variants} lists the command-line options of every evaluated variant.

\begin{table}[h]
\centering\footnotesize
\setlength{\tabcolsep}{3pt}
\begin{tabular}{@{}ll@{}}
\toprule
Variant & \texttt{run\_eval} options \\
\midrule
Released & \texttt{--method baseline} \\
Open loop & \texttt{--method open\_loop --k 5} (or 10) \\
Gripper events & \texttt{... --event-trigger --hold 5} \\
Stale image & \texttt{--method stale --k 5 --flow-steps 3} \\
Fewer steps & \texttt{--method flow --flow-steps 3} \\
Float16 & \texttt{... --dtype fp16} \\
CUDA graph & \texttt{... --graph} \\
Fixed noise & \texttt{... --fixed-noise} \\
\bottomrule
\end{tabular}
\caption{Evaluated variants.}
\label{tab:variants}
\end{table}

\section{Additional Results}
\label{app:results}

\paragraph{Per-task success.}
\cref{tab:pertask} gives the number of successful episodes out of 20 per task for each variant on initial states 0--19.

\begin{table}[h]
\centering\small
\setlength{\tabcolsep}{3.2pt}
\begin{tabular}{@{}lrrrrrrrrrr@{}}
\toprule
Task & 0 & 1 & 2 & 3 & 4 & 5 & 6 & 7 & 8 & 9 \\
\midrule
Released & 14 & 19 & 13 & 10 & 16 & 19 & 15 & 16 & 17 & 15 \\
bf16, $S{=}3$ & 15 & 18 & 15 & 14 & 15 & 18 & 14 & 15 & 19 & 15 \\
fp16, $S{=}3$ & 15 & 18 & 17 & 15 & 15 & 17 & 14 & 12 & 17 & 14 \\
Open loop $k{=}5$ & 12 & 14 & 13 & 11 & 13 & 13 & 12 & 11 & 16 & 14 \\
$k{=}5$ + events & 12 & 16 & 10 & 12 & 13 & 17 & 12 & 12 & 15 & 11 \\
Stale image & 11 & 13 & 17 & 12 & 13 & 16 & 11 & 13 & 12 & 12 \\
Open loop $k{=}10$ & 15 & 15 & 12 & 9 & 12 & 15 & 17 & 11 & 14 & 15 \\
\bottomrule
\end{tabular}
\caption{Successful episodes out of 20 per LIBERO-Spatial task.}
\label{tab:pertask}
\end{table}

\paragraph{Failure categories.}
Failures of the released policy are categorised from the executed gripper commands, since object poses are not logged: segments of at most 3 steps are first merged into their neighbour (chatter), and a closure is a switch from open to closed (\texttt{scripts/report\_numbers.py}).
\emph{Three or more closures}: 9 of 46.
\emph{Closed at the step limit} (one or two closures): 9.
\emph{Reopened after one or two closures}: 28.
Every failure closes the gripper at least once.

\paragraph{Warm-started flow (rejected before closed-loop tests).}
Starting the integration from the previous chunk, shifted by one step and re-noised to flow time $t_0$, was calibrated offline on 586 control steps of five LIBERO-Object episodes.
The deviation of the first action from a 10-step integration with the same noise was 0.83, 0.53 and 0.34 for $t_0=0.3, 0.5, 0.7$ (3 steps), against 0.32 for a cold start with 3 steps and 0.51 with 2 steps; the shifted previous chunk alone deviated by 1.31, and a 10-step integration with different noise by 1.37.
A warm start was therefore no more accurate than a cold start with the same number of steps and was not evaluated in closed loop.
